\section{Related Work}
\subsection{World Model}
Learning a dynamics model of the world that supports downstream action generation has been a long-standing challenge~\citep{ha2018world,sutton1991dyna}. Early work learns world models in low-dimensional state spaces~\citep{achille2018separation,lesort2018state}, which are efficient to train but difficult to generalize across visually diverse environments. With advances in generative modeling, a growing body of recent work has explored video models as foundation world models~\citep{kong2024hunyuanvideo,wang2025wan,yang2024cogvideox}. However, these models remain in the 2D pixel space, limiting their ability to capture 3D geometric structure and leaving a large representational gap between their predictions and the 3D actions a robot must produce.

3D world models attempt to close this gap by reasoning over meshes or explicit surfaces~\citep{wang2021neus,pfaff2020learning,jiang2025phystwin,zhao2025physsplat,xia2025drawer,xia2024video2game,zhen2025learning,guo2026articulat3d}, radiance fields or Gaussians~\citep{mildenhall2021nerf,kerbl20233d,driess2023learning,xie2024physgaussian}, or particle systems~\citep{sanchez2020learning,abou2024particlenerf,zhang2025particle,chen20254dnex}. Hybrid approaches additionally reason over hierarchical structures~\citep{kaelbling2011hierarchical,wang2025enact,zhao2026high}. While these representations offer richer geometric reasoning, they often require multi-view inputs, are scene-specific, or lack the scalability of pretrained video priors. Most closely related to our work are~\citep{zhen2025learning}, which models the 4D scene from RGB-DN (RGB, Depth, and Normal) videos with language-conditioned control, and~\citep{chen20254dnex}, which produces high-quality dynamic point clouds for novel-view video synthesis. 
Our work shares the goal of scalable 4D prediction but introduces a projective 4D representation that co-generates optical flow alongside RGB and depth, making inter-frame 3D motion explicit. Unlike methods relying on static geometry like surface normals~\citep{zhen2025learning}, our inclusion of optical flow allows for back-projection into 3D scene flow, providing explicit dynamic cues essential for learning accurate inverse dynamics. This is particularly critical for dexterous manipulation, where the fine-grained trajectory of objects and end-effectors is what differentiates success from failure.



\subsection{Future Prediction for Embodied Control}
A growing line of work bridges generative modeling and control by using 2D future prediction to guide policy learning~\citep{bharadhwaj2024gen2act,ye2024latent,ye2026world,bi2025motus,hu2024video}. Representative methods include SuSIE~\citep{black2023zero}, which employs a goal-conditioned keyframe generator~\citep{brooks2023instructpix2pix}, and UniPi~\citep{du2023learning}, which learns inverse dynamics over generated sequences. Downstream actions are then derived via online planning~\citep{hu2024video,williams2017model,hafner2019learning,pineau2003point}, offline policy synthesis~\citep{hafner2019dream,hansen2023td,chua2018deep,hafner2025training}, or inverse-dynamics models~\citep{du2023learning,bi2025motus}. However, because these pipelines operate entirely in 2D pixel space and often require repeated denoising for every action step, they face inherent limitations in both geometric accuracy and control reactivity.
In contrast, \textbf{\our} operates on a unified 4D representation that jointly encodes appearance, geometry, and motion. Building upon this, \textbf{\our-Policy} directly consumes the internal predictive features of \our in a single forward pass. By bypassing the need for per-step video decoding and denoising, our approach enables high-frequency, closed-loop robotic control, effectively translating imagined 4D trajectories into precise, real-time robotic actions.

